\documentclass[11pt]{article}
\usepackage[margin=1.1in]{geometry}
\usepackage{times}
\usepackage[T1]{fontenc}
\usepackage{setspace}
\onehalfspacing
\usepackage{parskip}
\pagenumbering{gobble}

\title{AGI by [Year]}
\author{Flyxion}
\date{September 2026}

\begin{document}
\maketitle

Artificial general intelligence will arrive in eighteen months. This estimate has held steady for the better part of a decade, which should be the headline finding and never is.

The prediction is not falsified when the date passes uneventfully. The date was based on the old scaling curve, the old benchmark suite, or an insufficiently generous definition of "general," and the revised forecast, incorporating everything learned from the miss, lands eighteen months out again. A weather service that missed this consistently would be replaced. A forecaster of elections who redefined "winning" every time their candidate lost would be laughed out of the trade. This forecast does neither, because it was never scored against the thing it named. It was scored against how confidently it was delivered, and confidence, unlike accuracy, does not decay on a schedule.

Notice the asymmetry doing the work. A hit is credited entirely to foresight. A miss is credited to nobody, because by the time it would be due for accounting the target has already moved and taken the record of the original claim with it. A method that wins on both outcomes has stopped measuring the world and started measuring the audience's patience for being told the same thing again with more urgency attached. The urgency is the tell. Genuine forecasts get calmer as they're refined by evidence. This one gets more dire with every miss, because the miss itself is repackaged as proof that things are accelerating, which is the one interpretation of a null result that no other forecasting discipline on earth would accept without asking what, specifically, accelerated.

\end{document}
